% Unnumbered insertion in A.77 before its final certificate paragraph.
\paragraph{The dyadic norm levels for the constructed pairs.}
The same fixed-level norm conclusion holds at $\ell=2$
for every ordinary pair under consideration with
$v_3(A)=1$ and $A/3$ a square in
$\widehat{\Q^\times}$.
In particular it applies to the point-$3$ congruent
ghosts constructed above.
Use the compatible rational representatives $a_n$
of A.77, with $m=2^n$.
Each $a_n$ has the same rational square class as $3$,
so it has the same $\Q_2$ square class as $3$.

The three quadratic subfields of
$\Q_2(\mu_{2^\infty})$ are
$\Q_2(i)$, $\Q_2(\sqrt2)$ and
$\Q_2(\sqrt{-2})$.
Indeed its Galois group is
$\Z_2^\times\simeq C_2\times\Z_2$; its three
index-two subgroups already occur in
$\Q_2(\mu_8)$.
The square class of $3$ differs from
$-1,2,-2$ in $\Q_2^\times/\Q_2^{\times2}$:
the last two have odd valuation, and
$3/(-1)\equiv5\pmod8$.
It follows that $a_n$ is not a square in
$\Q_2(\mu_m)$.
Kummer theory over this field therefore gives
\[
 [\Q_2(\mu_m,a_n^{1/m}):\Q_2(\mu_m)]=m .
\]
The cyclotomic extension at two itself has degree
$\varphi(m)$.
For the global splitting field
$L_n=\Q(\mu_m,a_n^{1/m})$ this proves
that the decomposition group at two has order
$m\varphi(m)$, the largest possible global order.
Consequently it is the whole group
$G_n=\operatorname{Gal}(L_n/\Q)$.

Here is the exact Hasse-principle implication, including
the inflation map.
For any number-field torus $T$ split by a finite Galois
extension $L/k$, let $X=X^*(T)$ and $G=\operatorname{Gal}(L/k)$.
Tate--Poitou duality identifies the norm-torus
Tate--Shafarevich group with the dual of the kernel of
\[
 H^2(G,X)\longrightarrow
                  \prod_v H^2(G_v,X).
\]
This is the general statement in
\cite[Proposition 6.4 and Corollary 6.5, p.~35]
{BIR_Oki2026Norm}; it is not restricted to the
prime-squared extensions studied elsewhere in that paper.
To verify the finite-group comparison directly, a class
in $\operatorname{Sha}^2(k,X)$ restricts over $L$ to a tuple of
characters, because
$H^2(G_L,\Z^r)=\operatorname{Hom}_{\rm cts}(G_L,\Q/\Z)^r$.
All these characters are locally trivial, hence vanish
by Chebotarev.
Moreover $H^1(G_L,X)=0$, since a continuous map from
a compact group to the discrete torsion-free lattice
$X$ has finite, therefore zero, image.
The Hochschild--Serre sequence thus puts the class
in the image of inflation from $H^2(G,X)$.
At a place with $G_v=G$, the same vanishing
$H^1(G_{L_w},X)=0$ makes local inflation injective.
Local triviality at that one place therefore kills
the inflated global class.
Duality proves $\operatorname{Sha}^1(k,T)=0$.
For $T=R^1_{K/k}\G_m$, Hilbert 90 identifies
$H^1(k,T)=k^\times/N_{K/k}K^\times$ and the local
maps with their ordinary norm maps.
This proves the Hasse norm principle for $K/k$.

Apply this argument to $K_n/\Q$ and its splitting
field $L_n$, with the full decomposition group at two
just proved.
The finite-place local norm identity (A.77.12) is
unchanged for $m=2^n$.
The rational number $a_n$ is positive, so $K_n$ has
a real embedding; the norm of its real etale algebra
onto $\R^\times$ is surjective.
Thus a rational representative $b_n$ of $B$ modulo
$m$ is a local norm at every place and is a global
norm.
Finally, the kernel of reduction modulo $2^n$ on
$\widehat{\Q^\times}_2$ is exactly its subgroup of
$2^n$th powers: the sign is already detected modulo
$2^n$, and the remaining prime coordinates can be
divided by $2^n$ with the same cofinal convergence.
The scalar correction $D_n$ in A.77 therefore gives
the actual completed norm preimage of $B$ also at two.
The remaining norm-tower compatibility class is
exactly the one already displayed in (A.77.13).
